% 编辑转录副本；不是独立下载原件。来源：https://github.com/bbadzioch/topology_lecture_notes/blob/master/MTH_727_notes/17_hurewicz_theorem.tex

Hurewicz homomorphism is a map that connects homotopy and homology groups.
Recall that $H_n(S^n)\cong\Z$. We will denote by $\gamma_n$ a chosen generator of $H_n(S^n)$. Given an element $[\varphi:(S^n,s_0)\to(X,x_0)]\in\pi_n(X,x_0)$ consider the homomorphism $\varphi_*:H_*(S^n)\to H_n(X)$. This homomorphism depends only on the homotopy class of $\varphi$.
\begin{definition}
The \emph{Hurewicz homomorphism} is a function
\[h:\pi_n(X,x_0)\to H_n(X)\]
given by $h([\varphi])=\varphi_*(\gamma_n)$.
\end{definition}
\begin{proposition}[Naturality]
For any function $f:X\to Y$ the following diagram commutes:
\[
\xymatrix{\pi_n(X,x_0)\ar[r]^{f_*}\ar[d]_h&\pi_n(Y,f(x_0))\ar[d]^h\\H_n(X)\ar[r]_{f_*}&H_n(Y)}
\]
\end{proposition}
\begin{proof}
For $[\varphi]\in\pi_n(X,x_0)$ we have
\[hf_*([\varphi])=h([f\varphi])=(f\varphi)_*(\gamma_n)=f_*\varphi_*(\gamma_n)=f_*h([\varphi]).\]
\end{proof}
\begin{proposition}
The Hurewicz homorphism is a group homomorphism.
\end{proposition}
\begin{proof}
Let $\varphi,\psi:(S^n,s_0)\to(X,x_0)$ where $n\geq1$. Recall that the element $[\varphi]\cdot[\psi]\in\pi_n(X,x_0)$ is the homotopy class of the map
\[S^n\overset p\longrightarrow S^n\vee S^n\overset{\varphi\vee\psi}\longrightarrow X\]
where $p:S^n\to S^n\vee S^n$ is the pinch map. Let $r_1,r_2:S^n\vee S^n\to S^n$ be the retractions of $S^n\vee S^n$ onto the first and, respectively, the second copy of $S^n$. We have a commutative diagram
\[
\xymatrix@C=3em{H_n(S^n)\ar[r]^{p_*}\ar[dr]_{\id_*\oplus\id_*}&H_n(S^n\vee S^n)\ar[r]^{(\varphi\vee\psi)_*}\ar[d]_{r_{1*}\oplus r_{2*}}^{\cong}&H_n(X)\\&H_n(S^n)\oplus H_n(S^n)\ar[ur]_{\varphi_*+\psi_*}&}
\]
This gives:
\begin{align*}
h([\varphi]\cdot[\psi])&=((\varphi\vee\psi)p)_*(\gamma_n)\\
&=(\varphi_*+\psi_*)(\id_*\oplus\id_*)(\gamma_n)\\
&=\varphi_*(\gamma_n)+\psi_*(\gamma_n)=h([\varphi])+h([\psi]).
\end{align*}
\end{proof}
\begin{theorem}[Hurewicz isomorphism theorem]
Let $X$ be a path connected space such that for some $n\geq2$ we have $\pi_i(X)=0$ for $i<n$. Then $H_i(X)=0$ for $0<i<n$ and the Hurewicz homomorphism
\[h:\pi_n(X,x_0)\to H_n(X)\]
is an isomorphism.
\end{theorem}
\begin{proof}
Assume first that $X=S^n$. We have $H_i(S^n)=0$ for $0<i<n$. In degree $n$ the Hurewicz homomorphism is a map $h:\Z\cong\pi_n(S^n)\to H_n(S^n)\cong\Z$. The group $\pi_n(S^n)$ is generated by the homotopy class of the identity map $\id_{S^n}:S^n\to S^n$ (\textsf{PIN SN NOTE}). We have $h([\id_{S^n}])=\id_{S^n*}(\gamma_n)=\gamma_n$. Therefore $h$ maps a generator of $\pi_n(S^n)$ to a generator of $H^n(S^n)$, and so it is an isomorphism.

Next, assume that $X=\bigvee_{i\in I}S^n$. Again, in this case $H_i(\bigvee_{i\in I}S^n)=0$ for $0<i<n$. Also, the retraction maps $r_i:\bigvee_{i\in I}S^n\to S^n$ give a commutative diagram
\[
\xymatrix{\pi_n(\bigvee_{i\in I}S^n)\ar[r]^{\bigoplus r_{i*}}_{\cong}\ar[d]_h&\bigoplus_{i\in I}\pi_n(S^n)\ar[d]^{\bigoplus_{i\in I}h}_{\cong}\\H_n(\bigvee_{i\in I}S^n)\ar[r]_{\bigoplus r_{i*}}^{\cong}&\bigoplus_{i\in I}H_n(S^n)}
\]
The the map $\bigoplus_{i\in I}h$ is an isomorphism by the previous case, so the left vertical map $h$ is also an isomorphism.

For the next step, assume that $X$ is an arbitrary CW complex with $\pi_i(X)=0$ for $i<n$. By Proposition \textsf{N SKELETON FOR N CONNECTED PROP} we can assume that $X^{(n-1)}=*$, which gives $H_i(X)=0$ for $0<i<n$.

Let $j:X^{(n+1)}\hookrightarrow X$ be the inclusion of the $(n+1)$-skeleton of $X$. By the naturality proposition we have a commutative diagram
\[
\xymatrix{\pi_n(X^{(n+1)})\ar[r]^{j_*}_{\cong}\ar[d]_h&\pi_n(X)\ar[d]^h\\H_n(X^{(n+1)})\ar[r]_{j_*}^{\cong}&H_n(X)}
\]
The upper homomorphism $j_*$ is an isomorphism by Proposition \textsf{PIN FOR CW SKELETON PROP}, and the lower $j_*$ is an isomorphism by properties of homology groups. As a consequence, it is enough to show that $h:\pi_n(X^{(n+1)})\to H_n(X^{(n+1)})$ is an isomorphism.

Since $X^{(n-1)}=*$, it follows that $X^{(n)}=\bigvee_{i\in I}S^n$ and $X^{(n+1)}=X^{(n)}\cup\bigcup_{k\in K}e_k^{(n+1)}$ where $\{e_k^{(n+1)}\}_{k\in K}$ are $(n+1)$-cells of $X$. Let $\varphi_k:S^n\to X^{(n)}$ be the attaching map of the cell $e_k^{n+1}$, and let $i:X^{(n)}\hookrightarrow X^{(n+1)}$ denote the inclusion map. We have a commutative diagram
\[
\xymatrix@C=1.1em{
\pi_n(\bigvee_{k\in K}S^n)\ar[r]^{(\bigvee\varphi_k)_*}\ar[d]_h^{\cong}&\pi_n(X^{(n)})\ar[r]^{i_*}\ar[d]_h^{\cong}&\pi_n(X^{(n+1)})\ar[r]\ar[d]^h&0\ar[d]^{\cong}\\
H_n(\bigvee_{k\in K}S^n)\ar[r]_{(\bigvee\varphi_k)_*}&H_n(X^{(n)})\ar[r]_{i_*}&H_n(X^{(n+1)})\ar[r]&H_{n-1}(\bigvee_{k\in K}S^n)
}
\]
The upper row of this diagram is exact by Proposition \textsf{1CONN PIN XCUPE PROP}, and the lower row is exact by the long homology sequence associated to the map $\bigvee_{k\in K}\varphi_k$. By the Five Lemma we obtain that $h:\pi_n(X^{(n+1)})\to H_n(X^{(n+1)})$ is an isomorphism.

Finally, let $X$ be an arbitary space with $\pi_i(X)=0$ for $i<n$. Let $f:Y\to X$ be a CW approximation of $X$ (\textsf{CW APPROX DEF}). Using Theorem \textsf{WE HOMOLOGY ISO THM} and the previous case we get $H_i(X)\cong H_i(Y)=0$ for $0<i<n$.

By the naturality proposition we have a commutative diagram
\[
\xymatrix{\pi_n(Y)\ar[r]^{f_*}_{\cong}\ar[d]_h^{\cong}&\pi_n(X)\ar[d]^h\\H_n(Y)\ar[r]_{f_*}^{\cong}&H_n(X)}
\]
Since $f$ is a weak equivalence, the upper homomorphism $f_*$ is an isomorphism by definition, and the lower $f_*$ is an isomorphism by Theorem \textsf{WE HOMOLOGY ISO THM}. Also, since $Y$ is a CW complex the left vertical map is an isomorphism by the previous case. Therefore $h:\pi_n(X)\to H_n(X)$ is an isomorphism.
\end{proof}
